\documentclass[11pt]{article}
\usepackage[margin=1in]{geometry}
\usepackage[T1]{fontenc}
\usepackage[utf8]{inputenc}
\usepackage{lmodern}
\usepackage{amsmath,amssymb,mathtools,booktabs,graphicx,natbib,hyperref}
\usepackage[nameinlink,noabbrev]{cleveref}

\title{Adaptive Forecast-Optimal Trend Estimation under Changing Time-Series Regimes}
\author{Heriberto Espino Montelongo}
\date{\today}

\newcommand{\R}{\mathbb{R}}
\newcommand{\Dt}{D_d}

\begin{document}
\maketitle

\begin{abstract}
We study forecast-optimal trend estimation as an adaptive forecasting method,
where smoothness, finite-memory window length, and finite-difference order may
depend on forecast horizon and local time-series regime. The central object is
therefore the full forecasting-method configuration
\(\Theta^\star_{T,h}=(d^\star_{T,h},L^\star_{T,h},S^\star_{T,h})\), not a
single penalty parameter. For the pure quadratic penalized trend, analytic
first- and second-order derivatives of future forecast loss with respect to the
penalty parameter permit derivative-aware numerical selection in log-penalty
space. The empirical design is strictly chronological. Controlled simulations
first separate forecast optimality from historical trend-recovery optimality,
then study mechanisms such as persistence and horizon, and finally examine
within-series regime transitions and adaptive-versus-fixed out-of-sample
performance. External evidence is planned from macroeconomic series through
market indices, equities, and cryptocurrency. This manuscript remains a
research draft; novelty and performance claims are deferred until the
literature audit and planned experiments are complete.
\end{abstract}

\section{Problem formulation}
Let $y\in\R^n$, let $\Dt$ be the order-$d$ finite-difference matrix, and let
$Q=\Dt^\top\Dt$. This finite-difference penalized-least-squares construction is closely related
to classical graduation/smoothing methods and to the controlled-smoothness
framework of \citet{guerrero2007penalized,whittaker1923graduation}. The pure
penalized trend used for the initial analytic development is
\begin{equation}
\label{eq:pure-trend}
\widehat t_{\lambda,d}
=
\arg\min_{t\in\R^n}
\left\{
\|y-t\|_2^2+\lambda\|\Dt t\|_2^2
\right\}
=
(I+\lambda Q)^{-1}y.
\end{equation}
Writing $S_\lambda=(I+\lambda Q)^{-1}$,
\begin{equation}
\label{eq:smoother-derivative}
S_\lambda'=-S_\lambda Q S_\lambda.
\end{equation}

\section{Forecast-optimal smoothness}
At forecast origin $T$, only observations available through $T$ may enter the
fit. Let $H$ denote the linear continuation operator from the fitted trend to
the next $h$ observations. Define
\begin{equation}
\label{eq:forecast-residual}
r_T(\lambda)
=
y_{T+1:T+h}
-
H S_\lambda y_{\mathrm{past}},
\end{equation}
and
\begin{equation}
\label{eq:forecast-loss}
f_T(\lambda)
=
\frac{1}{h}r_T(\lambda)^\top r_T(\lambda).
\end{equation}
The first derivative is
\begin{equation}
\label{eq:forecast-loss-first}
f_T'(\lambda)
=
\frac{2}{h}
r_T(\lambda)^\top
H S_\lambda Q S_\lambda y_{\mathrm{past}},
\end{equation}
and the second derivative is
\begin{equation}
\label{eq:forecast-loss-second}
f_T''(\lambda)
=
\frac{2}{h}
\left[
\|H S_\lambda Q S_\lambda y_{\mathrm{past}}\|_2^2
-
2r_T(\lambda)^\top
H S_\lambda Q S_\lambda Q S_\lambda y_{\mathrm{past}}
\right].
\end{equation}

For repeated chronological forecast origins $T_1,\ldots,T_M$, the active
validation objective pools future-block squared errors:
\begin{equation}
\label{eq:rolling-cv}
CV_{d,L,h}(\lambda)
=
\frac{\sum_{j=1}^M h_j f_{T_j}(\lambda)}
{\sum_{j=1}^M h_j}.
\end{equation}
The first and second derivatives are obtained by the same weighted aggregation
of \cref{eq:forecast-loss-first,eq:forecast-loss-second}.

Because the normalized smoothness index depends on sample length, the principal
inner-validation design uses a fixed-width window $L$ for each candidate
configuration. Thus all inner fits for a given $(d,L)$ use the same sample size,
and a common $\lambda$ corresponds to one common normalized smoothness level.
When several values of $L$ are compared, they are scored on the same set of
inner forecast origins so that changing $L$ does not also change the validation
targets.

\subsection{Forecasting versus reconstruction}
Randomly omitting interior observations can be useful when the scientific
target is smoothing, interpolation, or trend reconstruction. It is not the
validation design used for the forecasting claims in this paper, because
observations later than an interior omission are then available to the
smoother. Here, at every origin $T$, model fitting and hyperparameter selection
use only information available through $T$; the future block is revealed only
after the forecast has been formed.

\section{Numerical selection}
We work with
\begin{equation}
\label{eq:log-lambda}
\theta=\log\lambda,
\qquad
g(\theta)=CV(e^\theta).
\end{equation}
Then
\begin{equation}
\label{eq:log-derivatives}
g'(\theta)=\lambda CV'(\lambda),
\qquad
g''(\theta)=\lambda CV'(\lambda)+\lambda^2CV''(\lambda).
\end{equation}
The planned robust procedure uses a coarse log-penalty scan only to bracket
sign changes of $g'$, applies a bracketed root solver to every detected
interval using a bracketed zero-finding method in the spirit of
\citet{brent1971zero}, classifies the stationary points, and compares all local
minima with the domain boundaries. Dense grid search and Newton iterations are
retained as diagnostic and computational benchmarks.

\section{Research questions}
At forecast origin \(T\) and horizon \(h\), define the forecast-method
configuration
\begin{equation}
\label{eq:adaptive-configuration}
\Theta^\star_{T,h}
=
\left(
d^\star_{T,h},
L^\star_{T,h},
S^\star_{T,h}
\right)
=
G\!\left(h,X_T,\mathcal C\right),
\end{equation}
where \(d\) is difference order, \(L\) is finite-memory estimation-window
length, \(S\) is normalized smoothness, \(X_T\) summarizes the local state, and
\(\mathcal C\) denotes the series class. Candidate components of \(X_T\)
include volatility, persistence, trend strength, local roughness, and break
evidence.

The paper has three empirical questions. First, does the forecast-optimal
configuration differ from the configuration that best reconstructs a known
historical latent trend? Second, which mechanisms make
\(\Theta^\star_{T,h}\) vary with horizon and local state? Third, does adapting
the forecasting method to those states improve untouched future forecasts
relative to strong fixed configurations?

The persistence/horizon experiment is one mechanism study addressing the
second question. Its oracle AR-residual-aware forecast is diagnostic only: it
must not be interpreted as the definition of the paper or as the proposed
real-data model. Local regime is deliberately broader than persistence alone.

\section{Experimental design}
The planned sequence is:
\begin{enumerate}
\item controlled simulations separating forecast and recovery targets;
\item mechanism studies for persistence, horizon, signal-to-noise structure,
roughness, endpoints, and related effects;
\item within-series regime-transition experiments tracking the joint path of
\((d^\star,L^\star,S^\star)\) and adaptation delay;
\item adaptive-versus-fixed untouched out-of-sample comparisons;
\item lower-frequency macroeconomic series;
\item broad market indices and ETFs;
\item individual equities;
\item cryptocurrency.
\end{enumerate}
All real-data experiments use nested chronological evaluation. At an outer
origin, only the available history is passed to the inner selector; after
$(d,L,\lambda)$ is selected, the chosen model is refitted on the selected
window and forecasts the untouched future block. For financial level forecasts,
the mandatory baseline is
\begin{equation}
\label{eq:no-change}
\widehat y^{(0)}_{T+k\mid T}=y_T,
\qquad k=1,\ldots,h,
\end{equation}
and we report the relative pooled forecast error
\begin{equation}
\label{eq:relative-rmsfe}
RMSFE_{\mathrm{rel}}
=
\frac{RMSFE_{\mathrm{method}}}
{RMSFE_{\mathrm{no-change}}}.
\end{equation}
Return or trend-change targets will be considered separately.

\section{Scope}
The current paper does not study portfolio utility or decision-optimal
smoothing. Those questions remain possible future work built on the same
\texttt{trend\_estimation} library.

\bibliographystyle{plainnat}
\IfFileExists{literature/references.bib}{
  \bibliography{literature/references}
}{
  \bibliography{../literature/references}
}

\end{document}
